\chapter[Troubleshooting]{Typical Throwing and Construction Problems}
\label{sec:DebugThrowing}
\begin{quote}
THE LESSON\\
\\
I threw it out, it went up high.\\
It came back down right at my eye.\\
I dove to the ground with some alarm,\\
Never again will I throw it side arm.\\\\
\footnotesize \textit{Bud Pell}
\end{quote}
\newpage
\section{Flies straight ahead}
\begin{center}
	\includegraphics[width=30em]{2STRAT.eps} 
\end{center}
\begin{itemize}
\item Either a right-hander has thrown a left-handed boomerang or the other way around. 
\item The profiled side (top) of the boomerang was not facing towards the thrower when launched
\item The boomerang is a non-returner, also called a hunting stick.
\item The boomerang has too much negative dihedral (see section \ref{sec:dihedral}) 
\item The boomerang has too little airfoil - not enough lift. This is often true of commercial boomerangs. 
\item The boomerang is too heavy. 
\item There was too little spin when launched. 
\item The boomerang has too many holes or flaps.
\item You twisted your wrist instead of throwing it straight.
\end{itemize}
\newpage

\section[Sharp rise, then crash]{Flies toward the ground, flies sharply upward and then dives down out of control}
\begin{center}
	\includegraphics[width=30em]{2FLAT.eps} 
\end{center}
\begin{itemize}
\item The tilt was almost 90° (sidearm throw), when it should have been approximately 20° from vertical (see section \ref{sec:Throwing}). 
\item The boomerang could be warped. Normally the wing tips and the rest of the boomerang should lie flat on a counter. 
\item Arm 1 has too much lift. 
\item The average maximum lift of the boomerang is too far ahead of the centre of rotation (see chapter \ref{sec:ConstructionDisk}).
\end{itemize}
\newpage

\section[Circular but touches ground]{Flight path is circular but the boomerang touches the ground during flight}
\begin{center}
	\includegraphics[width=30em]{2VERT.eps} 
\end{center}
\begin{itemize}
\item The throw was too weak.
\item The boomerang was thrown too high.
\item The tilt is too small or even negative.
\item At least one wing has a negative dihedral (see chapter \ref{sec:ConstructionDisk}).
\item The shape of the boomerang has an average maximum lift behind the centre of rotation (see chapter \ref{sec:ConstructionDisk}).
\item Arm 2 has too much lift compared to Arm 1.
\item If this happens to a multiblader, the hole in the middle is too big.
\end{itemize}
\newpage

\section[Crosses the thrower]{Crosses in front of the thrower and lands to the right (RH) / left (LH)}
\begin{center}
	\includegraphics[width=25em]{2RIGHT.eps} 
\end{center}
\begin{itemize}
\item The angle to the wind is too big (too far ``out of the wind'' see section \ref{sec:Throwing}). 
\end{itemize}
\newpage

\section[Passes behind the thrower]{Passes behind the thrower and lands to the left (RH) / right (LH)}
\begin{center}
	\includegraphics[width=25em]{2LEFT.eps} 
\end{center}
\begin{itemize}
	\item The angle to the wind is too small (too close ``into the wind'' see section \ref{sec:Throwing}).
\end{itemize}
\newpage

\section[Returns almost]{Returns toward the thrower, but falls short}
\begin{center}
	\includegraphics[width=30em]{2TILT.eps} 
\end{center}
\begin{itemize}
\item Throw lacked sufficient power.
\item There was too little spin when launched.
\item The tilt is too small or even negative.
\item The boomerang has too many holes or flaps or is too rough.
\item Arm 2 has too little lift compared to arm 1.
\item The multiblader has too little lift (see chapter \ref{sec:correct}). 
\item Check the wind direction. There is no wind. The boomerang needs more wind.
\end{itemize}
\newpage

\section{Returns over and beyond the thrower}
\begin{center}
	\includegraphics[width=30em]{2HARD-A.eps} 
\end{center}
\begin{itemize}
\item Thrown with not enough tilt (see section \ref{sec:Throwing}). 
\item The throw was too low (see section \ref{sec:Throwing}).
\item Arm 1 has too little lift compared to arm 2.
\item The multiblader has too much lift (see chapter \ref{sec:correct}). 
\item Check the wind. The wind is too strong.
\end{itemize}
\newpage

\section[Flies upward then crashes]{Flies upward during the first half of the flight but then crashes during the second half}
\begin{center}
	\includegraphics[width=30em]{2HIGH.eps} 
\end{center}
\begin{itemize}
\item The height of throw is too great (thrown too high see section \ref{sec:Throwing}).
\end{itemize}	
\newpage

\section[Close to ground then rises]{Flies close to the ground during the first half of the flight path but then rises during the second half.}
\begin{center}
	\includegraphics[width=30em]{2LOW.eps} 
\end{center}
\begin{itemize}
\item The height of throw is too small (thrown too low see section \ref{sec:Throwing}).
\item Arm 1 has too much lift compared to arm 2
\end{itemize}
\newpage

\section[No layover]{Does not lay-over - maybe even rolls inward}
\begin{center}
	\includegraphics[width=30em]{2IN.eps} 
\end{center}
\begin{itemize}
\item Arm 2 has too much lift in relation to arm 1. (This sometimes depends on the shape of the boomerang, i.e. if arm 2 is bigger, wider and thicker than arm 1).
\item The shape of  the boomerang has average maximum lift behind the centre of rotation (see chapter \ref{sec:ConstructionDisk})
\item The boomerang is warped. One or more arms have too much negative dihedral, or Arm 1 has negative angle of attack (see p. 14).
\end{itemize}
\newpage

\section[S-shaped flight]{Flight path of the boomerang looks like an „S“ in the final stage}
\begin{center}
	\includegraphics[width=30em]{2S.eps} 
\end{center}
\begin{itemize}
\item Arm 1 has too much lift compared to Arm 2. 
\item Because of its shape the boomerang has average maximum lift too far ahead of the centre of rotation (see chapter \ref{sec:ConstructionDisk})
\item The boomerang is warped.
\item The height of throw is too small (thrown too low see section \ref{sec:Throwing}).
\end{itemize}
\newpage

\section[Extra loops]{Flies a few extra loops in the final stage}
\begin{center}
	\includegraphics[width=30em]{2HARD.eps} 
\end{center}
\begin{itemize}
\item Thrown too hard! 
\item The boomerang has too little drag (see chapter \ref{sec:correct}).
\end{itemize}
\newpage

\section{Flies perfectly}
\begin{center}
	\includegraphics[width=30em]{2GOOD.eps} 
\end{center}
\begin{itemize}
\item Great job! You got it.
\end{itemize}
Experienced throwers often purposely use the above-mentioned ``errors'', to make use of their effects during certain situations (e.g.competitions or dealing with unusually windy or calm conditions). Correctness of a throw is subjective. A seemingly erroneous flight may be perfect for the conditions. A throw is ``good'' if it achieves the return the thrower wants or needs.
